\documentclass[11pt]{article}
\usepackage[T1]{fontenc}
\usepackage{textcomp}
\usepackage[margin=0.75in]{geometry}
\usepackage{amsmath,amssymb,amsthm}
\usepackage{array,booktabs}
\usepackage{hyperref}
\setlength{\parindent}{0pt}
\newtheorem{theorem}{Theorem}
\theoremstyle{definition}
\newtheorem{definition}{Definition}
\title{Flow Proof Artifact: Eq-symmetric}
\author{Generated by \texttt{flow doc proof}}
\date{Flow Proof Book}
\begin{document}
\maketitle
If two naturals are equal, the equality reads either way.\\[0.5em]
\textbf{Source.} leibniz — https://en.wikipedia.org/wiki/Identity\_(mathematics)\#Equality\\[0.5em]
\bigskip
\noindent{\large \textbf{Derived fact 1.} \textit{Equality reverses: from x = y we obtain y = x} \hfill equality \cdot equality \cdot equality reverses}\par
\smallskip\noindent\textit{Coordinate: equality \cdot equality \cdot equality reverses}\par
\smallskip\noindent\textit{Source: leibniz — https://en.wikipedia.org/wiki/Symmetric\_relation}\par
\smallskip\noindent\textit{Built on: anything is always equal to itself}\par
\medskip
\noindent\textbf{Goal.} Equality reverses: from x = y we obtain y = x\par
\noindent$\displaystyle \forall x \in \mathbb{N} \forall y \in \mathbb{N}\quad y = x$\par
\medskip
\noindent
\begin{tabular}{@{} >{\raggedright\arraybackslash}p{0.06\textwidth} >{\raggedright\arraybackslash}p{0.47\textwidth} >{\raggedleft\arraybackslash}p{0.41\textwidth} @{}}
\textbf{\#} & \textbf{Proof} & \textbf{Mathematics} \\
\midrule
\textbf{1.} & We prove that equality reverses for equality on equality. & \\[0.45em]
\textbf{2.} & We split into exhaustive cases — the claim must hold in each one. & \\[0.45em]
\textbf{3.} & Case 1 (see step 2): suppose x  equals  y. & \\[0.45em]
\textbf{4.} & We invoke the axiom governing equality on equality: anything is always equal to itself (instantiated for x). & $\displaystyle x = x$ \\[0.45em]
\textbf{5.} & We invoke the axiom governing equality on equality: anything is always equal to itself (instantiated for y). & $\displaystyle y = y$ \\[0.45em]
\textbf{6.} & From step 3, step 4, and step 5, this implies y equals x. Together with the other cases (step 3), the goal is discharged. Hence proven. & $\displaystyle y = x$ \\[0.45em]
\bottomrule
\end{tabular}
\label{thm:Eq-equality-equality_reverses}
\par\medskip\hrule
\end{document}
